\documentclass[10pt,a4paper]{amsart}
\usepackage[margin=19mm]{geometry}
\usepackage{amsmath,amssymb,mathtools}
\usepackage[scr=rsfs,cal=euler]{mathalfa}
\usepackage{xfrac}
\usepackage[unicode,hidelinks,bookmarks=false]{hyperref}
\newcommand{\ie}{i.e.\ }
\newcommand{\cf}{cf.\ }
\newcommand{\resp}{respectively\ }
\newcommand{\Subsection}{\subsection}
\providecommand{\nobreakdash}{\nobreak\hbox{-}}
\setlength{\emergencystretch}{2em}
\multlinegap0pt
\usepackage{bm}
\usepackage{bbm}
\usepackage{dsfont}
\usepackage{xspace}
\usepackage{cancel}
\usepackage{mathtools}
\usepackage{amsthm}
\makeatletter
\@ifundefined{Theorem}{\newtheorem{Theorem}{Theorem}}{}
\@ifundefined{Proposition}{\newtheorem{Proposition}[Theorem]{Proposition}}{}
\makeatother
\DeclareMathOperator{\CC}{\mathbb{C}}
\DeclareMathOperator{\OO}{\mathcal{O}}
\DeclareMathOperator{\WW}{{\bf W}}
\DeclareMathOperator{\XX}{{\bf X}}
\DeclareMathOperator{\YY}{{\bf Y}}
\title[The Ising Model Coupled to 2D Gravity: Genus Zero Partition ]{The Ising Model Coupled to 2D Gravity: Genus Zero Partition Function}
\author{Maurice Duits, Nathan Hayford, Seung-Yeop Lee}
\date{Source: arXiv:2405.03259v3 (CC BY 4.0)}
\begin{document}
\maketitle

\noindent\emph{Source and licence.} The Ising Model Coupled to 2D Gravity: Genus Zero Partition Function. Version arXiv:2405.03259v3 (CC BY 4.0), distributed under the Creative Commons Attribution 4.0 International licence, as recorded in the arXiv licence metadata for this exact version and shown on the abstract page https://arxiv.org/abs/2405.03259v3. Full rights notice: licenses/arxiv-2405.03259-CC-BY-4.0.txt.\par\smallskip

\noindent\textit{Editorial scope.} Counted unit: the Proposition whose source label is \texttt{X-asymptotics} in the article (source0.tex lines 1977--2020, source label \texttt{X-asymptotics}), delivered together with the article's own complete original proof of it (source0.tex lines 2022--2112, one \texttt{proof} environment of 91 lines). No mathematics is written, repaired or re-derived: statement and proof are the article's own text line for line. Required context retained uncounted: Lambda-Matrix (lines 1788--1854). Every definition, display and label of those lines is retained unchanged. Omitted from this package and not redistributed: 1--1787, 1855--1976, 2021--2021, 2113--4952. Nothing inside a retained excerpt is omitted or altered: every retained body is delivered exactly as the article's lines are written, so no omission receipt of any kind accompanies this package. Citations: the retained excerpts carry no citation command at all, so no bibliography entry of the article is transcribed and no reference is redistributed. Reference adaptation: none was needed and none was made; every reference inside the delivered unit resolves inside it, so no unresolved reference or citation is produced. Printed slips kept literal: no printed word, symbol, hypothesis or number of the counted statement or of its proof was altered. Layout and class: the article's own class is \texttt{sigma} with packages including ifpdf, mathtools, xcolor, tikz, pgfplots; the wrapper is the lane's pinned \texttt{amsart} editorial wrapper at 10pt on A4, which restates the theorem-like environments the retained text needs and adds the article's own macro definitions that the retained text needs (\texttt{\textbackslash CC}, \texttt{\textbackslash OO}, \texttt{\textbackslash WW}, \texttt{\textbackslash XX}, \texttt{\textbackslash YY}), with extra wrapper packages (bm, bbm, dsfont, xspace, cancel, mathtools). The delivered numbering is the unit's own. Assets: no retained span contains a figure environment, a graphics-inclusion command, a PDF-page inclusion, a table or a TikZ picture; no image file is copied into this package, the package declares no asset dependency and no image file is redistributed with it. Licence: version arXiv:2405.03259v3 of this article is distributed under the Creative Commons Attribution 4.0 International licence, as recorded in the arXiv licence metadata for this exact version and shown on the abstract page https://arxiv.org/abs/2405.03259v3. A full rights notice is delivered at licenses/arxiv-2405.03259-CC-BY-4.0.txt. Characters: every retained excerpt is pure ASCII, so the delivered engine, XeLaTeX with its default Latin Modern text font, reports no missing character.
 \begin{Proposition} $\WW(z)$ is the unique solution to the following Riemann--Hilbert problem:
 \begin{enumerate}\itemsep=0pt
 \item[$(1)$] $\WW(z)$ is analytic in $\CC \setminus (\Gamma_1\cup \Gamma_2)$, with boundary values
 \begin{equation*}
 \WW_{+}(z) =
 \begin{cases}
 \begin{pmatrix}
 1 & 0 & 1 \\
 0 & 1 & 0 \\
 0 & 0 & 1
 \end{pmatrix}\WW_{-}(z) =: J_1\WW_{-}(z), & z \in \Gamma_1,\\[1.5mm]
 \begin{pmatrix}
 1 & 1 & 0 \\
 0 & 1 & 0 \\
 0 & 0 & 1
 \end{pmatrix}\WW_{-}(z)=: J_2\WW_{-}(z), & z \in \Gamma_2
 \end{cases}
 \end{equation*}
 $($note that these matrices commute, and that their product is $J_1J_2=J)$.
 \item[$(2)$] At infinity, $\WW(z)$ is normalized as
 \begin{equation*}
 \WW(z) = c_N \cdot {\rm e}^{N\Theta(z)} A(z) B(z)\hat{K}\biggl[\mathbb{I}_{3\times 3} + \OO\biggl(\frac{1}{z}\biggr)\biggr], \qquad z\to \infty,
 \end{equation*}
 where the constant $c_N:= {\rm i}\sqrt{\frac{2\pi}{3N}}{\rm e}^{N\frac{q}{6t}}$, the matrix $\Theta(z)$ is defined to be
 \begin{equation} \label{Lambda-Matrix}
 \Theta(z) =
 \begin{cases}
 {\rm diag}(\lambda_3(z), \lambda_1(z),\lambda_2(z)), & \operatorname{Im} z > 0,\\
 {\rm diag}(\lambda_2(z), \lambda_1(z),\lambda_3(z)), & \operatorname{Im} z <0,
 \end{cases} %{\rm diag}(\theta_{1}(z),\theta_{2}(z),\theta_{3}(z))
 \end{equation}
where the functions $\lambda_k(z)$ are defined by the \textit{exact} formulas
 \begin{equation*}%\label{little-theta-hat}
 \lambda_k(z) = -\frac{3\omega^{k-1}}{4} \frac{\tau^{4/3}}{(-tq^{-1})^{1/3}}z^{4/3} - \frac{\omega^{1-k}}{2}\frac{\tau^{2/3}}{(-tq^{-1})^{2/3}}z^{2/3}.
 \end{equation*}
and finally, the matrices $A(z)$, $B(z)$, and $\hat{K}$ are given by
 \begin{gather*}
 A(z) =
 \begin{cases}
 \begin{pmatrix}
 -\omega & 1 & \omega^2\\
 -1 & 1 & 1\\
 -\omega^2 & 1 & \omega
 \end{pmatrix}, & \operatorname{Im} z > 0,\\[2.5mm]
 \begin{pmatrix}
 \omega^2 & -1 & -\omega\\
 -1 & \hphantom{-}1 & \hphantom{-}1\\
 -\omega & \hphantom{-}1 & \hphantom{-}\omega^2\\
 \end{pmatrix}, & \operatorname{Im} z < 0,
 \end{cases}
 \qquad
 B(z) =
 \begin{pmatrix}
 z^{-1/3} & 0 & 0\\
 0 & 1 & 0\\
 0 & 0 & z^{1/3}
 \end{pmatrix},
 \\
 \hat{K} =
 \begin{pmatrix}
 \bigl(-tq^{-1}\bigr)^{-1/6}\tau^{-1/3} & 0 & -\frac{n+27 tq^{-1}}{54(-tq^{-1})^{13/6}\tau^{1/3}}\\[1.4mm]
 0 & \bigl(-tq^{-1}\bigr)^{-1/2} & 0\\[1.1mm]
 0 & 0 & -\bigl(-tq^{-1}\bigr)^{-5/6}\tau^{1/3}
 \end{pmatrix}.
 \end{gather*}
 \end{enumerate}
 \end{Proposition}

 \begin{Proposition}
 $\XX(z)$ is the unique solution to the following Riemann--Hilbert Problem:
 \begin{equation*}
 \XX_{+}(z) = \XX_{-}(z) \times
 \begin{cases}
 \begin{pmatrix}
 1 & 0\\
 0 & J_1^{-1}
 \end{pmatrix}, & z \in \Gamma_1,\vspace{1mm}\\
 \begin{pmatrix}
 1 & 0\\
 0 & J_2^{-1}
 \end{pmatrix}, & z \in \Gamma_2,\vspace{1mm}\\
 \begin{pmatrix}
 1 & 1 \ 0 \ 0\\
 \vec{0}_{1\times 3} & \mathbb{I}_{3\times3}
 \end{pmatrix}, & z \in \Gamma,\\
 \end{cases}
 \end{equation*}

 Subject to the normalization condition
 \begin{align}\label{X-asymptotics}
 \XX(z) =
 \biggl[\mathbb{I} + \OO\biggl(\frac{1}{z}\biggr)\biggr]\begin{pmatrix}
 1 & 0\\
 0 & B^{-1}(z) A^{-1}(z)
 \end{pmatrix}
 \begin{pmatrix}
 z^n & 0 & 0 & 0\\
 0 & z^{-n/3} & 0 & 0\\
 0 & 0 & z^{-n/3} & 0\\
 0 & 0 & 0 & z^{-n/3}
 \end{pmatrix}{\rm e}^{-N\Lambda(z)},
 \end{align}
where $\Lambda(z)$ is defined as the diagonal matrix
 \begin{equation*}
 \Lambda(z) =
 \begin{pmatrix}
 V(z) & 0\\
 0 & \Theta(z)
 \end{pmatrix},
 \end{equation*}
where $\Theta(z)$ is the diagonal $3\times 3$ matrix defined by \eqref{Lambda-Matrix}.
 \end{Proposition}

\begin{proof}
 The asymptotic condition \eqref{X-asymptotics} follows almost immediately from the definition of $\XX(z)$; along with the definition of the matrix $\WW(z)$. Indeed, this is obvious in the regions $\Omega_u$ and $\Omega_\ell$, and follows from definition of $\WW(z)$. The only detail to check
 is that the asymptotics of $\WW(z)$ in the region $\Omega_c$ are correct. We see that $\WW(z)$ in this region is obtained by adding the recessive solution $-\vec{w}_1(z)$
 to the first row; since this solution is recessive in the region $\Omega_c$, the asymptotics of $-\vec{w}_2(z) - \vec{w}_1(z)$ $(= \vec{w}_3(z)+ \vec{w}_4(z))$ are the same as the
 asymptotics of $-\vec{w}_2(z)$ there. Thus, the asymptotics of $\XX(z)$ from the
 proposition hold.

 Now, we address the jump conditions of $\XX(z)$. Since the matrix function
 \begin{equation*}
 \begin{pmatrix}
 {\rm e}^{-NV(z)} & 0 \\
 0 & \WW^{-1}(z)
 \end{pmatrix}
 \end{equation*}
 has jumps only on $\Gamma_1$, $\Gamma_2$ (arising from the jumps of $\WW(z)$), and these contours do not intersect~$\Gamma$, the first two jump conditions are clearly satisfied. It remains to check the jump of $\XX(z)$ across~$\Gamma$. The jump of $\YY(z)$ on $\Gamma$ is
 \begin{equation*}
 J_{\YY} = \mathbb{I} +
 {\rm e}^{-NV(z)}\begin{pmatrix}
 0 & f(z) & \frac{f'(z)}{N\tau} & \frac{f''(z)}{(N\tau)^2} \\
 0 & 0 & 0 & 0 \\
 0 & 0 & 0 & 0 \\
 0 & 0 & 0 & 0
 \end{pmatrix}
 =:
 \mathbb{I} + {\rm e}^{-NV(z)} \begin{pmatrix}
 0 & & \vec{f}(z) & \\
 0 & 0 & 0 & 0 \\
 0 & 0 & 0 & 0 \\
 0 & 0 & 0 & 0
 \end{pmatrix},
 \end{equation*}
 where $\vec{f}(z)$ is the row vector
 \begin{equation*}
 \vec{f}(z) := \left(
 f(z), \frac{f'(z)}{N\tau}, \frac{f''(z)}{(N\tau)^2} \right).
 \end{equation*}
 Now, since $\Gamma$ is homologically equivalent to $-\gamma_2 - \gamma_1$, we have that $\vec{L}(z) = -\vec{w}_2(z) - \vec{w}_1(z)$. Thus, the jump of $\XX(z)$ across $\Gamma$ is
 \begin{gather*}
 \begin{pmatrix}
 {\rm e}^{NV(z)} & 0 \\
 0 & \WW(z)
 \end{pmatrix}
 \left[
 \mathbb{I} + {\rm e}^{-NV(z)}
 \begin{pmatrix}
 0 & & \vec{f}(z) & \\
 0 & 0 & 0 & 0 \\
 0 & 0 & 0 & 0 \\
 0 & 0 & 0 & 0
 \end{pmatrix}\right]
 \begin{pmatrix}
 {\rm e}^{-NV(z)} & 0 \\
 0 & \WW^{-1}(z)
 \end{pmatrix}\\
 \qquad{}= \mathbb{I} +
 \begin{pmatrix}
 0 & & \vec{f}(z)\WW^{-1}(z) & \\
 0 & 0 & 0 & 0 \\
 0 & 0 & 0 & 0 \\
 0 & 0 & 0 & 0
 \end{pmatrix} = \mathbb{I} +
 \begin{pmatrix}
 0 & & (-\vec{w}_2(z) - \vec{w}_1(z))\WW^{-1}(z) & \\
 0 & 0 & 0 & 0 \\
 0 & 0 & 0 & 0 \\
 0 & 0 & 0 & 0
 \end{pmatrix} \\
 \qquad{}=
 \begin{pmatrix}
 1 & 1 & 0 & 0\\
 0 & 1 & 0 & 0 \\
 0 & 0 & 1 & 0 \\
 0 & 0 & 0 & 1
 \end{pmatrix}.
 \end{gather*}
 Here we have used the relations
 \begin{alignat*}{3}
 &\vec{w}_1(z)\WW^{-1}(z) = (0,0,1), \qquad&& \vec{w}_2(z)\WW^{-1}(z) = (-1,0,-1),&\\
 &\vec{w}_3(z)\WW^{-1}(z) = (0,1,0), \qquad&& \vec{w}_4(z)\WW^{-1}(z) = (1,-1,0),&
 \end{alignat*}
 resulting from the fact that, in a neighborhood of $\Gamma$, the matrix $\WW(z)$ admits the expression
 \begin{equation*}
 \WW(z) =
 \begin{pmatrix}
 \vec{w}_3(z) + \vec{w}_4(z) \\
 \vec{w}_3(z) \\
 \vec{w}_1(z)
 \end{pmatrix},
 \end{equation*}
 along with the identity $\WW(z) \WW^{-1}(z) = \mathbb{I}_{3\times 3}$.
\end{proof}

\end{document}
